\documentclass[10pt]{article}
\usepackage[T1]{fontenc}
\usepackage{lmodern,amsmath,amssymb,amsthm,microtype}
\usepackage[margin=0.82in]{geometry}
\usepackage[colorlinks=true,urlcolor=blue]{hyperref}
\newtheorem{theorem}{Theorem}
\newtheorem{lemma}{Lemma}
\newtheorem{corollary}{Corollary}
\title{Records and Unit Values of an XOR-Triangle Boundary}
\author{Qichao Wang\\{\small Hebei University of Technology}\\[0.3em]
Daoyu Dong\\{\small University of Electronic Science and Technology of China}}
\date{October 26, 2026}
\begin{document}
\maketitle
\begin{abstract}
We prove both conjectures in OEIS A334595: its record positions are
one and $2^m+1$ for $m\ge1$, and its value equals one exactly at
powers of two. The right boundary of an XOR triangle is a
binomial transform over the field with two elements. Its inverse
is the same transform, which determines every value, every
fixed-length image, and the unique maximum of each length block.
\end{abstract}

\section{The boundary transform}
For $n\ge1$, form the XOR triangle whose first row is the ordinary
binary expansion of $n$. Each entry in a subsequent row is the
sum modulo two of the two entries above it. Read the right boundary
from the upper right corner to the bottom as a binary word;
its integer value is $a(n)$ in OEIS A334595.

Write $n=\sum_{i=0}^{L-1}b_i2^i$, with $b_{L-1}=1$, and let
$y_j$ be the rightmost entry in row $j$, starting with $j=0$.
Repeated addition of neighboring entries gives
\begin{equation}\label{eq:transform}
 y_j=\sum_{i=0}^{j}\binom ji b_i\pmod2,
 \qquad
 a(n)=\sum_{j=0}^{L-1}y_j2^{L-1-j}.
\end{equation}
Pascal's identity proves the first expression inductively.
This is a finite explicit all-index formula, independent of
constructing the interior of the triangle.

\begin{lemma}\label{lem:inverse}
The map $T:(b_i)\mapsto(y_j)$ in \eqref{eq:transform} is its own
inverse on every fixed length $L$ over $\mathbf F_2$.
\end{lemma}
\begin{proof}
For $0\le i\le j<L$, the coefficient of $b_i$ in $(T^2b)_j$ is
\[
 \sum_{k=i}^j\binom jk\binom ki
 =\binom ji\sum_{k=i}^j\binom{j-i}{k-i}
 =\binom ji2^{j-i}.
\]
Modulo two this equals one for $j=i$ and zero for $j>i$.
\end{proof}

\section{All values equal to one}
\begin{theorem}
For every positive integer $n$, $a(n)=1$ if and only if $n$ is a
power of two.
\end{theorem}
\begin{proof}
On length $L$, output one has boundary vector
$(y_0,\ldots,y_{L-1})=(0,\ldots,0,1)$.
Its inverse under $T$ is that same vector: for $j<L-1$ every
input summand is zero, and at $j=L-1$ the diagonal entry is one.
Thus $b_{L-1}=1$ is its only nonzero input bit, giving
$n=2^{L-1}$. Conversely this input clearly gives output one.
\end{proof}

\section{Every record position and value}
\begin{theorem}\label{thm:records}
For every $L\ge2$, the maximum of $a(n)$ on
$2^{L-1}\le n<2^L$ is $2^L-2$, attained uniquely at
$n=2^{L-1}+1$. Therefore the strict record positions of A334595
are exactly
\[
 1,\quad 2^m+1\quad(m\ge1),
\]
with respective record values $1$ and $2^{m+1}-2$.
\end{theorem}
\begin{proof}
All outputs are represented by $L$ bits, so at most $2^L-1$.
The all-one boundary vector has inverse with
\[
 b_0=1,\qquad b_j=\sum_{i=0}^j\binom ji=2^j=0\pmod2
 \quad(j\ge1).
\]
For $L\ge2$ its leading input bit is zero; hence output $2^L-1$
is impossible on this input length. The next possible output,
$2^L-2$, has vector $(1,\ldots,1,0)$. It differs from the
all-one vector only in its last coordinate, so its inverse has
exactly the first and last coordinates equal to one.
This is the valid input $n=2^{L-1}+1$, unique by
Lemma~\ref{lem:inverse}. It is therefore the unique block maximum.

The first input in each block is $2^{L-1}$ and has output one
by the preceding theorem. The second input already attains the
block maximum, which exceeds the maxima of all previous length
blocks. No later input in that block can set a new strict record.
The initial block consists of $n=1$ with output one.
\end{proof}

The fixed-length image itself also has an exact description:
it consists of the $L$-bit words $y$, allowing leading zeros,
that satisfy
\[
 \sum_{i=0}^{L-1}\binom{L-1}{i}y_i=1\pmod2.
\]
This is precisely the condition that their unique inverse has
leading bit one. Thus no exceptional word or index is hidden in
the characterization above.

\section{Source and independent checks}
The entry's linked MathOverflow discussion concerns fixed points
of a related boundary rotation. Gerry Myerson's May 2, 2020
comment observes the binomial nature of XOR differences. We use
that standard transform explicitly and prove the two specific
conjectures still stated in A334595; no priority for the transform
itself is claimed.
The verifier compares all 8191 official terms with direct triangle
construction and an independent binomial-coefficient implementation.
It checks both conjectures and the unique block maxima through
$n=65535$, and verifies $T^2=I$ on every binary vector of lengths
one through ten. The all-index conclusions follow from the proofs.

\begin{thebibliography}{2}
\bibitem{oeis} The OEIS Foundation Inc.,
\href{https://oeis.org/A334595}{A334595} and its
\href{https://oeis.org/A334595/b334595.txt}{8191-term b-file},
accessed October 26, 2026.
\bibitem{mo} DSM and contributors,\ \emph{Number triangle},
\url{https://mathoverflow.net/questions/359138/number-triangle},
May 2020, accessed October 26, 2026.
\end{thebibliography}
\end{document}
